% Document/LinearAlgebra/VectorSpaces/Matrices/Part5_Invertibility.tex
% Reordered Chapter 3 -- Section 5 (closes the Matrices chapter): invertibility of matrices (as maps),
% ([T]_{B_U}^{B_V})^{-1} = [T^{-1}]_{B_V}^{B_U}; T invertible iff its matrix is; every change-of-basis
% matrix invertible. Equivalence + similarity were SPLIT OUT into EquivalenceSimilarity.tex (a separate
% chapter). General cardinalities. DRAFT for review.

\newpage
\section{Invertibility}
\label{sec:invertibility}

\begin{intuition}
    With maps and matrices now identified (Section~\ref{sec:matrices-as-maps}), ``inverting a matrix''
    means exactly ``inverting the map it is,'' and the arithmetic of Section~\ref{sec:matrices} guarantees
    the inverse is again a matrix. One question then organises this section: how does inversion interact
    with the matrix of an \emph{abstract} map? The answer is as clean as possible --- inverting the map
    swaps the roles of the two bases and inverts the matrix --- and it specialises to the fact that every
    change-of-basis matrix is invertible, its inverse being the change run backwards. (Which matrices
    represent one and the same map, the question this naturally raises next, is taken up in the following
    chapter, on equivalence and similarity.)
\end{intuition}

\subsection{Invertible Matrices}

\begin{definition}[Invertible matrix]
\label{def:invertible-matrix}
    A matrix $A \in \CFMat{Y}{X}[K]$ is \textbf{invertible} if it is invertible as a linear map
    $A : K^{(X)} \to K^{(Y)}$ --- that is, bijective. Its inverse map is again linear, hence a matrix
    $A^{-1} \in \CFMat{X}{Y}[K]$ (the grid of $A^{-1}$ under the identification of
    Definition~\ref{def:identification}), characterised by
    \[
        A^{-1} A = \IdentityMat{X}, \qquad A A^{-1} = \IdentityMat{Y} .
    \]
\end{definition}

\begin{remark}[Two-sided inverse, and when one can exist]
\label{rem:invertible-square}
    Because the identification turns composition into the matrix product and the identity map into the
    identity grid (Corollary~\ref{cor:inherited-laws}), $A$ is invertible exactly when it has a two-sided
    multiplicative inverse $B$ with $B A = \IdentityMat{X}$ and $A B = \IdentityMat{Y}$, and then
    $B = A^{-1}$ is unique (a two-sided inverse in a monoid is unique). An invertible $A$ is an
    isomorphism $K^{(X)} \xrightarrow{\sim} K^{(Y)}$, so $\Card{X} = \Card{Y}$ (isomorphic coordinate
    spaces have canonical bases of equal cardinality); in the finite case $X = Y = n$ and $A$ is a
    \emph{square} $n \times n$ matrix. The invertible elements of the algebra
    $\End[K]{K^{(X)}} = \CFMat{X}{X}[K]$ form, under the matrix product, the group $\GL{K^{(X)}}$ ---
    written $\GL{K^n}$ when $X = n$.
\end{remark}

\subsection{The Inverse of the Matrix of a Map}

\begin{theorem}[Inversion swaps the two bases]
\label{thm:inverse-matrix}
    Let $T : U \to V$ be an isomorphism, with ordered bases $\mathcal B_U$ of $U$ and $\mathcal B_V$ of
    $V$. Then $\CoordMatrix[\mathcal B_U][\mathcal B_V]{T}$ is invertible and
    \[
        \bigl( \CoordMatrix[\mathcal B_U][\mathcal B_V]{T} \bigr)^{-1}
        = \CoordMatrix[\mathcal B_V][\mathcal B_U]{T^{-1}} .
    \]
\end{theorem}

\begin{proof}
    Since $T$ is an isomorphism, $T^{-1} : V \to U$ is linear, so
    $\CoordMatrix[\mathcal B_V][\mathcal B_U]{T^{-1}}$ is defined. By the composition theorem
    (Theorem~\ref{thm:composition-matrix}) and the fact that the matrix of an identity in one basis is
    the identity grid (Remark~\ref{rem:change-of-basis}, with $\Phi(\Identity{}) = I$ from
    Corollary~\ref{cor:inherited-laws}),
    \[
        \Composition{\CoordMatrix[\mathcal B_V][\mathcal B_U]{T^{-1}}}{\CoordMatrix[\mathcal B_U][\mathcal B_V]{T}}
        = \CoordMatrix[\mathcal B_U][\mathcal B_U]{\Composition{T^{-1}}{T}}
        = \CoordMatrix[\mathcal B_U][\mathcal B_U]{\Identity{U}}
        = \IdentityMat{X} ,
    \]
    and symmetrically
    $\Composition{\CoordMatrix[\mathcal B_U][\mathcal B_V]{T}}{\CoordMatrix[\mathcal B_V][\mathcal B_U]{T^{-1}}}
    = \CoordMatrix[\mathcal B_V][\mathcal B_V]{\Composition{T}{T^{-1}}}
    = \CoordMatrix[\mathcal B_V][\mathcal B_V]{\Identity{V}} = \IdentityMat{Y}$. By
    Remark~\ref{rem:invertible-square} the two-sided inverse is unique, which is the claim.
\end{proof}

\begin{corollary}[Invertibility is detected by any matrix]
\label{cor:invertible-detect}
    For $T : U \to V$ with ordered bases $\mathcal B_U, \mathcal B_V$, the map $T$ is an isomorphism iff
    $\CoordMatrix[\mathcal B_U][\mathcal B_V]{T}$ is invertible. In particular every change-of-basis
    matrix is invertible, with
    \[
        \bigl( \CoordMatrix[\mathcal B][\mathcal C]{\Identity{V}} \bigr)^{-1}
        = \CoordMatrix[\mathcal C][\mathcal B]{\Identity{V}} .
    \]
\end{corollary}

\begin{proof}
    If $T$ is an isomorphism, Theorem~\ref{thm:inverse-matrix} exhibits the inverse. Conversely,
    $\CoordMatrix[\mathcal B_U][\mathcal B_V]{T} = \Composition{\Chart{\mathcal B_V}}{\Composition{T}{\Frame{\mathcal B_U}}}$
    with $\Chart{\mathcal B_V}, \Frame{\mathcal B_U}$ isomorphisms (Definition~\ref{def:matrix-of-map}); if
    this composite is bijective then
    $T = \Composition{\Frame{\mathcal B_V}}{\Composition{\CoordMatrix[\mathcal B_U][\mathcal B_V]{T}}{\Chart{\mathcal B_U}}}$
    is a composite of bijections, hence bijective. The change-of-basis matrix is the case
    $T = \Identity{V}$, an isomorphism, and its inverse is read off from
    Theorem~\ref{thm:inverse-matrix}.
\end{proof}

\subsection{Worked Example}

\begin{example}[The inverse of a matrix is the matrix of the inverse]
\label{ex:inverse-shear}
    Let $T : \bb{R}^2 \to \bb{R}^2$ be the shear $T(x, y) = (x + y,\; y)$, an isomorphism with inverse
    $T^{-1}(x, y) = (x - y,\; y)$. In the standard basis $\mathcal B$,
    \[
        \CoordMatrix[\mathcal B][\mathcal B]{T} = \begin{pmatrix} 1 & 1 \\ 0 & 1 \end{pmatrix},
        \qquad
        \CoordMatrix[\mathcal B][\mathcal B]{T^{-1}} = \begin{pmatrix} 1 & -1 \\ 0 & 1 \end{pmatrix},
    \]
    and indeed
    $\begin{psmallmatrix} 1 & 1 \\ 0 & 1 \end{psmallmatrix}\begin{psmallmatrix} 1 & -1 \\ 0 & 1 \end{psmallmatrix}
    = \begin{psmallmatrix} 1 & 0 \\ 0 & 1 \end{psmallmatrix}$, confirming
    $\bigl(\CoordMatrix[\mathcal B][\mathcal B]{T}\bigr)^{-1} = \CoordMatrix[\mathcal B][\mathcal B]{T^{-1}}$
    as Theorem~\ref{thm:inverse-matrix} predicts.
\end{example}
